\chapter{Architettura}

\section{Stack}

\begin{itemize}[leftmargin=*]
  \item \textbf{Backend}: Python 3.11+, FastAPI, SQLAlchemy 2.0,
        Pydantic 2, uvicorn, OR-Tools, NumPy, scikit-learn,
        Faker.
  \item \textbf{Frontend}: SvelteKit (Svelte 4), Vite 5, Tailwind CSS,
        JavaScript puro.
  \item \textbf{DB}: SQLite via SQLAlchemy. Migrazioni idempotenti
        nel codice (no Alembic).
  \item \textbf{Engine I/O}: pickle Python; \texttt{engine\_io.py}
        converte fra DB e dict per il solver.
\end{itemize}

\section{Diagramma a blocchi}

\begin{figure}[h]
\centering
\begin{tikzpicture}[node distance=0.6cm and 0.6cm,
  every node/.style={draw, rounded corners=4pt, very thick,
                     minimum width=4.5cm, minimum height=1cm,
                     align=center, font=\small}]
  \node[fill=cBox!30] (fe) {SvelteKit / Vite\\\textit{16 tabs}};
  \node[fill=cBox!30, below=of fe] (be) {FastAPI backend\\\textit{18 router, 118 endpoint}};
  \node[fill=cBox!30, below=of be] (db) {SQLite\\\textit{$\sim$30 tabelle}};
  \node[fill=cBox!30, below=of db] (en) {CP-SAT engine\\\textit{decomp. spettrale + meta}};
  \draw[-{Stealth[length=3mm]}, thick] (fe) -- node[draw=none, right]
        {\footnotesize \texttt{/api/*} via Vite proxy} (be);
  \draw[-{Stealth[length=3mm]}, thick] (be) -- node[draw=none, right]
        {\footnotesize SQLAlchemy 2.0} (db);
  \draw[-{Stealth[length=3mm]}, thick] (be) to[bend left=45]
        node[draw=none, right]
        {\footnotesize \texttt{engine\_io.py} (pickle)} (en);
\end{tikzpicture}
\caption{Architettura a blocchi della webapp.}
\label{fig:arch}
\end{figure}

\section{Struttura cartelle}

\begin{lstlisting}[style=shell]
timetable/
+-- README.md                    -- landing page
+-- docs/                        -- manuale tecnico
+-- experiments/                 -- solver code + pickle snapshots
|   +-- cpsat_v2_assignment.py   -- Phase A
|   +-- cpsat_v2_timetable.py    -- Phase B
|   +-- decomposition_spectral_v2.py
|   +-- metaheuristics.py        -- LNS / SA / TS / ILS
|   +-- classroom_assignment.py  -- Phase 4
+-- schedule/                    -- legacy single-file prototypes
+-- webui/
|   +-- start.bat                -- launcher Windows
|   +-- start.sh                 -- launcher Linux/macOS
|   +-- stop.sh                  -- stop helper
|   +-- backend/                 -- FastAPI app
|   +-- frontend/                -- SvelteKit
|   +-- data/timetable.db        -- SQLite
|   +-- docs/                    -- guide brevi user-facing
+-- proposals/                   -- design notes + benchmarks
+-- *.pdf                        -- reference papers
\end{lstlisting}

\section{Lifecycle del backend}

\texttt{webui/backend/main.py} definisce un \texttt{lifespan} che
chiama \texttt{init\_db()} allo startup. Esso fa:

\begin{enumerate}
  \item \texttt{Base.metadata.create\_all(bind=engine)} -- crea le
        tabelle mancanti.
  \item \texttt{\_apply\_lightweight\_migrations()} -- ALTER TABLE
        idempotenti per i campi aggiunti in versioni successive
        (esempio: \texttt{school\_classes.curriculum\_id},
        \texttt{teachers.last\_name/first\_name/nickname}, ecc.).
        Ogni ALTER \`e guardato da \texttt{PRAGMA table\_info} per non
        duplicare.
\end{enumerate}

\section{Architettura del solver: dal DSL a CP-SAT}
\label{sec:arch-dsl-cpsat}

Step 1--4 del piano multi-day hanno introdotto un'architettura
\emph{a strati uniforme} per tutti gli usi del solver. Il
diagramma sotto riassume i flussi di trasformazione.

\begin{verbatim}
+-----------------------------------------------------------+
| UI (SvelteKit)                                            |
|   - Tab Vincoli: matrici + DSL + Plessi                   |
|   - Tab Workflow: pipeline (mode, granularity)            |
+-----------------+----------------------------+------------+
                  |                              |
        POST /api/constraints/general     POST /api/optimize/*
                  |                              |
                  v                              v
+-----------------+----------------+   +---------+---------+
| dsl_parser (general_dsl.py)      |   | router            |
|   parse() -> AST                  |   | run_async         |
+-----------------+----------------+   +---------+---------+
                  |                              |
                  v                              v
+-----------------+----------------+   +---------+---------+
| dsl_translator                   |   | engine pipeline   |
|   tabelle special-purpose        |   | (mono / decomp /  |
|   -> stringhe DSL canoniche       |   |  CG / BP)         |
+-----------------+----------------+   +---------+---------+
                  |                              |
                  +------------- merged ---------+
                                  |
                                  v
+--------------------------------+--------------------------+
| dsl_to_cpsat.DSLConstraintCompiler                         |
|   AST + slot dict + plessi_data                            |
|   -> add_assertion() / add_implication() su CpModel        |
+--------------------------------+--------------------------+
                                 |
              +------------------+------------------+
              v                  v                  v
       MonolithicSolver   PhaseBDaySolver   BP Pricer (9 grani)
       (cpsat_v2_*) +     (legacy +         + RF nodi
       seed via DSL       DSL augment)
              |                  |                  |
              +------------------+------------------+
                                 |
                                 v
                         soluzione (Lessons)
                                 |
                                 v
                          DSL evaluator post-hoc
                          (HARD: feasibility check)
                          (SOFT: contributo allo
                                  score globale)
\end{verbatim}

I punti chiave:
\begin{itemize}
\item \textbf{Un parser, un AST}. Tutte le sorgenti (matrici,
  DSL UI, regole plesso, vincoli legacy seedati) vengono
  unificate in stringhe DSL canoniche e parsate da
  \texttt{general\_dsl.py}.
\item \textbf{Un compilatore}. Il
  \texttt{DSLConstraintCompiler} \`e la porta unica verso
  CP-SAT. Mono solver, pricer di BP, PhaseBDaySolver, nodi
  Ryan-Foster: tutti applicano lo stesso compilatore al loro
  modello, garantendo coerenza dei vincoli fra le diverse
  pipeline.
\item \textbf{Zero-drift}.
  \texttt{seed\_implicit\_hardcoded(profs)} traduce in DSL i
  vincoli HARD legacy hardcoded. Le tuple di slot proibite
  emesse via DSL coincidono bit-per-bit con quelle del path
  legacy: la migrazione \`e progressiva e reversibile.
\item \textbf{SOFT post-hoc, HARD up-front}. I HARD del DSL
  diventano vincoli CP-SAT prima della solve. I SOFT del DSL
  sono attualmente valutati \emph{post-hoc}: il loro costo
  contribuisce allo score globale ma non orienta ancora la
  ricerca CP-SAT (la wired-up del SOFT in obiettivo \`e
  pianificata nei prossimi step).
\end{itemize}

% =========== Cap 3: Modello dati ===========
